\documentclass[10pt,a4paper]{amsart}
\usepackage[margin=19mm]{geometry}
\usepackage{amsmath,amssymb,mathtools}
\usepackage[scr=rsfs,cal=euler]{mathalfa}
\usepackage{xfrac}
\usepackage[unicode,hidelinks,bookmarks=false]{hyperref}
\newcommand{\ie}{i.e.\ }
\newcommand{\cf}{cf.\ }
\newcommand{\resp}{respectively\ }
\newcommand{\Subsection}{\subsection}
\providecommand{\nobreakdash}{\nobreak\hbox{-}}
\setlength{\emergencystretch}{2em}
\multlinegap0pt
\usepackage{bm}
\usepackage{bbm}
\usepackage{dsfont}
\usepackage{xspace}
\usepackage{cancel}
\usepackage{mathtools}
\usepackage{amsthm}
\makeatletter
\@ifundefined{theorem}{\newtheorem{theorem}{Theorem}}{}
\@ifundefined{lemma}{\newtheorem{lemma}[theorem]{Lemma}}{}
\@ifundefined{proposition}{\newtheorem{proposition}[theorem]{Proposition}}{}
\makeatother
\newcommand\Aut{\mathrm{Aut}}
\newcommand\Q{\mathcal{Q}}
\newcommand\aug{\fboxsep=-\fboxrule\!\!\!\fbox{\strut}\!\!\!}
\renewcommand\P{\mathcal{P}}
\renewcommand\S{\mathrm{S}}
\title[Regular bipartite multigraphs have many (but not too many) s]{Regular bipartite multigraphs have many (but not too many) symmetries}
\author{Peter J. Cameron, Coen del Valle, Colva M. Roney-Dougal}
\date{Source: arXiv:2405.20002v2 (CC BY 4.0)}
\begin{document}
\maketitle

\noindent\emph{Source and licence.} Regular bipartite multigraphs have many (but not too many) symmetries. Version arXiv:2405.20002v2 (CC BY 4.0), distributed under the Creative Commons Attribution 4.0 International licence, as recorded in the arXiv licence metadata for this exact version and shown on the abstract page https://arxiv.org/abs/2405.20002v2. Full rights notice: licenses/arxiv-2405.20002-CC-BY-4.0.txt.\par\smallskip

\noindent\textit{Editorial scope.} Counted unit: the Proposition whose source label is \texttt{1} in the article (source0.tex lines 390--392, source label \texttt{1}), delivered together with the article's own complete original proof of it (source0.tex lines 393--423, one \texttt{proof} environment of 31 lines). No mathematics is written, repaired or re-derived: statement and proof are the article's own text line for line. Required context retained uncounted: top (lines 218--221); minN (lines 239--241); centbound (lines 279--281); bestmat1 (lines 381--383). Every definition, display and label of those lines is retained unchanged. Omitted from this package and not redistributed: 1--217, 222--238, 242--278, 282--380, 384--389, 424--747. Nothing inside a retained excerpt is omitted or altered: every retained body is delivered exactly as the article's lines are written, so no omission receipt of any kind accompanies this package. Citations: the retained excerpts carry no citation command at all, so no bibliography entry of the article is transcribed and no reference is redistributed. Reference adaptation: none was needed and none was made; every reference inside the delivered unit resolves inside it, so no unresolved reference or citation is produced. Printed slips kept literal: no printed word, symbol, hypothesis or number of the counted statement or of its proof was altered. Layout and class: the article's own class is \texttt{daj} with packages including amsmath, latexsym, amssymb, amsthm; the wrapper is the lane's pinned \texttt{amsart} editorial wrapper at 10pt on A4, which restates the theorem-like environments the retained text needs and adds the article's own macro definitions that the retained text needs (\texttt{\textbackslash Aut}, \texttt{\textbackslash P}, \texttt{\textbackslash Q}, \texttt{\textbackslash S}, \texttt{\textbackslash aug}), with extra wrapper packages (bm, bbm, dsfont, xspace, cancel, mathtools). The delivered numbering is the unit's own. Assets: no retained span contains a figure environment, a graphics-inclusion command, a PDF-page inclusion, a table or a TikZ picture; no image file is copied into this package, the package declares no asset dependency and no image file is redistributed with it. Licence: version arXiv:2405.20002v2 of this article is distributed under the Creative Commons Attribution 4.0 International licence, as recorded in the arXiv licence metadata for this exact version and shown on the abstract page https://arxiv.org/abs/2405.20002v2. A full rights notice is delivered at licenses/arxiv-2405.20002-CC-BY-4.0.txt. Characters: every retained excerpt is pure ASCII, so the delivered engine, XeLaTeX with its default Latin Modern text font, reports no missing character.
\begin{lemma}\label{top}
Let $(\mathcal{P}_1,\mathcal{P}_2)$ be a pair of $(k,l)$-partitions with $N=M(\P_1,\P_2)$, let $K=\S_{k}\times \S_{k}$, and set $G=\S_{k\times l}$. 
Then the pointwise stabiliser $H:=G_{\mathcal{P}_1,\mathcal{P}_2}$ is permutation isomorphic to $$\left(\prod_{i,j\in [k]}\mathrm{Sym}(P_{1i}\cap P_{2j})\right) :  K_{N}.$$ In particular, $|H|=|K_{N}|\prod_{i,j}n_{ij}!$, and $H\cong G_{\Q_1,\Q_2}$ whenever $M(\Q_1,\Q_2)=N$. Moreover, $\P_1$ and $\P_2$ are both partwise fixed by $H$ if and only if $K_N=1$.
\end{lemma}

\begin{lemma}\label{minN}
Let $s$ and $t$ be positive integers and let $X$ be a multiplicity sequence. Set $x=x(s,t,X)$, $b=b(s,t,X)$, and $S=S(s,t,X)$. Then $\min\left\{\prod_{a\in A} a! : A\in S\right\}$ is achieved uniquely by the integer multiset $B\in S$ containing each rank with its associated multiplicity, and all other elements as close to each other as possible. Thus, this minimum is precisely $$\left(\prod_{i}r_i!^{m_i}\right)\lceil x\rceil!^b\lfloor x\rfloor!^{t-\left(\sum m_i\right)-b}.$$ 
\end{lemma}

\begin{lemma}\label{centbound}
Let $k\geq3$ and $q$ be positive integers, and $G=\S_{k\times kq}$. Let $N=(n_{ij})_{k\times k}$ be a $(k,qk)$-intersection matrix with largest entry $q+1$, occurring exactly once in each row and column. Then $|K_N|\geq 3$ and  $|G(N)|\geq (q+1)!^kq!^{k^2-2k}(q-1)!^k\cdot 3$.
\end{lemma}

\begin{lemma}\label{bestmat1}
Suppose $r=\pm 1$. Let $$N=\begin{cases} B(k,r)&\text{ if $k$ is odd}\\C(k,r)&\text{ if $k$ is even.}\end{cases}$$ Then $N$ is a partwise fixed $(k,l)$-intersection matrix.
\end{lemma}

\begin{proposition}\label{1}
Suppose $r=\pm1$. Then $$\chi(k,l)=(q-r)!^{\lceil k/2\rceil}q!^{k^2-k-2\lceil k/2\rceil}(q+r)!^{k+\lceil k/2\rceil}.$$ Moreover, if $\Gamma$ is a $(k,l)$-bipartite graph with $|\Aut(\Gamma)|=\chi(k,l)$, then either $\Aut_V(\Gamma)=1$ and $$\Gamma^*=\{(q-r)^{\lceil k/2\rceil},(q)^{k^2-k-2\lceil k/2\rceil},(q+r)^{k+\lceil k/2\rceil}\},$$ or $l=k+r$, $|\Aut_V(\Gamma)|=2$, and $$\Gamma^*=\{(1-r)^{\lceil k/2\rceil-1},(1)^{k^2-k-2\lceil k/2\rceil+2},(1+r)^{k+\lceil k/2\rceil-1}\}.$$
\end{proposition}

\begin{proof}
Let $N$ be as in Lemma~\ref{bestmat1}. From Lemma~\ref{top}, we deduce that \begin{equation}\label{G(N)1}|G(N)|=(q-r)!^{\lceil k/2\rceil}q!^{k^2-k-2\lceil k/2\rceil}(q+r)!^{k+\lceil k/2\rceil}.\end{equation} 
Let $A=(a_{ij})$ be a $(k,qk+r)$-intersection matrix with $|G(A)|\leq |G(N)|$. 

If $r=1$ then let $m$ be the number of entries of $A$ less than $q$, and if $r=-1$ then let $m$ be the number of entries of $A$ greater than $q$. If $m\geq\lceil k/2\rceil$, then by \eqref{G(N)1} and Lemma~\ref{minN} applied to $(s,t,X)=(lk,k^2,[(\lceil k/2\rceil,q-r)])$, since $$ x(s,t,X)=\frac{lk-(\lceil k/2\rceil)(q-r)}{k^2-\lceil k/2\rceil}$$ is between $q$ and $q+r$ we deduce that $|G(A)|=|G(N)|$, that $A$ is partwise fixed and that $A^*=N^*$. Therefore, we may assume $m<\lceil k/2\rceil$.

By reordering the rows and columns so that the $m$ distinguished entries occur in the initial $m\times m$ submatrix, $A$ has the form
$$\begin{pmatrix}
A_1&\aug&A_2\\\hline
A_3&\aug&A_4\\
\end{pmatrix}$$
where $A_1$ is $m\times m$ and $A_2,A_3,A_4$ have all entries in $\{q,q+r\}$  (for if not then they would also contain one of the distinguished entries, contradicting that there are only $m$ such entries). In particular, since all row and column sums are $qk+r$ we deduce that $q+r$ occurs exactly once in $R_i$ and $C_i$ for each $i>m$. For $1\leq i\leq m$, let $\mathcal{C}_i$ 
 be the set of all $j\in \{m+1,m+2,\dots, k\}$ such that $a_{ij}=q+r$ --- these entries appear in $A_2$ --- and similarly $\mathcal{R}_i$ 
 the set of all $ j\in \{m+1,m+2,\dots, k\}$ such that $a_{ji}=q+r$ --- these entries appear in~$A_3$. 

If $j_1,j_2\in \mathcal{C}_i$ for $1\leq i\leq m$, then $(1, (j_1\, j_2))\in K_A$, and similarly if $j_1,j_2\in \mathcal{R}_i$ then $((j_1\, j_2),1)\in K_A$. Finally, if $j_1,j_2\in \mathcal{C}_0$, then there exist unique $i_1\ne i_2>m$ such that $a_{i_1j_1}=a_{i_2j_2}=q+r$ and so $((i_1\, i_2), (j_1\, j_2))\in K_A$. Combining this with Lemma~\ref{minN} applied to $(s,t,X)=(lk,k^2,[(m,q-r)])$ gives \begin{equation}\label{asize}|G(A)|\geq (q-r)!^mq!^{k^2-k-2m}(q+r)!^{k+m}|\mathcal{C}_0|!\prod_{i=1}^{m}|\mathcal{C}_i|!|\mathcal{R}_i|!.\end{equation} We now split into cases according to the value of $m$.

\vspace{11pt}\noindent \underline{\textbf{Case $m<\lceil k/2\rceil -1$:}} The matrix $A_2$ has $k-m$ columns and at most one $q+r$ per column, so at most $m+1$ columns are distinct. Therefore, $|\mathcal{C}_0|!\prod_{i=1}^{m}|\mathcal{C}_i|!|\mathcal{R}_i|!\geq 2^{k-2m-1}$, so by \eqref{asize} $$|G(A)|\geq 2^{k-2m-1}(q-r)!^mq!^{k^2-k-2m}(q+r)!^{k+m}=\frac{2^{k-2m-1}q^{\lceil k/2\rceil-m}}{(q+1)^{\lceil k/2\rceil-m}}|G(N)|>|G(N)|,$$ a contradiction.

\vspace{11pt}\noindent \underline{\textbf{Case $m=\lceil k/2\rceil -1$:}} If two columns of $A_2$ or two rows of $A_3$ are the same, then by \eqref{asize}, $$|G(A)|\geq2(q-r)!^mq!^{k^2-k-2m}(q+r)!^{k+m}=\frac{2q}{q+1}|G(N)|\geq |G(N)|,$$ with equality only possible if $q=1$; $|K_A|=2$; and $m$ of the entries are equal to $(1-r)$, $k+m$ entries to $(1+r)$, and all other entries $1$. 
 On the other hand if all columns of $A_2$ are distinct, as well as all rows of $A_3$, then each row of $A_2$ has sum $(k-m)q+1$, as does each column of $A_3$, hence $A_1$ is an $(m,qm)$-intersection matrix. If $A$ is not partwise fixed then as before $|G(A)|\geq\frac{2q}{q+1}|G(N)|.$ Therefore, there remain two possibilities for $A$:
\begin{itemize}
    \item[(i)] $q=1$, $A_1$ is an $(m,m)$-intersection matrix with $m$ entries equal to $1-r$, all other entries in $\{1+r,1\}$,  and $|K_A|=2$.
    \item[(ii)] $A$ is partwise fixed with $A_1$ a partwise fixed $(m,qm)$-intersection matrix, all columns of $A_2$ are distinct, and all rows of $A_3$ are distinct.

\end{itemize}
In the first case we are done, so consider case (ii) --- we show it is not possible. If the $m$ distinguished entries of $A_1$ are not all in distinct rows and columns, then without loss of generality $R_{m}(A)$ begins with $(q)^m$. Since $m=\lceil k/2\rceil-1$ and all rows of $A_3$ are distinct, some $R_i(A)$ with $m+1\leq i\leq k$ begins $(q)^m$. Let $j_1,j_2\geq m+1$ be the unique indices such that $a_{mj_1}=a_{ij_2}=q+r$. Then $((m\, i), (j_1\,j_2))\in K_A$, and so $|K_A|\geq2$, a contradiction, whence the distinguished entries of $A_1$ are in distinct rows and columns.

Suppose $A_1$ has an entry greater than $q+1$ if $r=-1$ or less than $q-1$ if $r=1$. Then from Lemma~\ref{minN} we deduce that $$|G(A)|\geq (q-2r)!(q-r)!^{m-1}q!^{k^2-k-2m-1}(q+r)!^{k+m+1}=\frac{q-r+1}{q-r}|G(N)|>|G(N)|,$$
a contradiction, hence the $m$ distinguished entries all equal $q-r$. Since each $q-r$ occurs in a distinct row and column, $q+sr\not\in A_1^*$ for any $s>1$ --- that is, $A_1$ is as in Lemma~\ref{centbound}, so is not partwise fixed, a contradiction.
\end{proof}

\end{document}
